\chapter{Comparison of Gaussian Classifier Variants}
\label{chap:variants_comparison}

This chapter expands our analysis to compare three distinct Gaussian classifier variants: the Standard MVG, the Tied Covariance MVG, and the Naive Bayes classifier. We evaluate these models across three specific applications represented purely in terms of their effective priors ($\tilde{\pi} = 0.1$, $\tilde{\pi} = 0.5$, and $\tilde{\pi} = 0.9$) with uniform costs.

Table \ref{tab:comparison_results} summarizes the Actual $\text{DCF}_n$ and Minimum DCF extracted for these configurations.

\begin{table}[htbp]
    \centering
    \small 
    \begin{tabular}{llccc}
        \toprule
        \textbf{Eff. Prior ($\tilde{\pi}$)} & \textbf{Model} & \textbf{Actual $\text{DCF}_n$} & \textbf{Min DCF} & \textbf{Calib. Loss (\%)} \\
        \midrule
        \multirow{3}{*}{\textbf{0.5}} 
        & Standard MVG & 0.140 & 0.130 & 7.7\% \\
        & Naive Bayes  & 0.144 & 0.131 & 9.9\% \\
        & Tied MVG     & 0.186 & 0.181 & 2.7\% \\
        \midrule
        \multirow{3}{*}{\textbf{0.9}} 
        & Standard MVG & 0.400 & 0.342 & 17.0\% \\
        & Naive Bayes  & 0.389 & 0.351 & 10.8\% \\
        & Tied MVG     & 0.463 & 0.442 & 4.7\% \\
        \midrule
        \multirow{3}{*}{\textbf{0.1}} 
        & Standard MVG & 0.305 & 0.263 & 16.0\% \\
        & Naive Bayes  & 0.302 & 0.257 & 17.5\% \\
        & Tied MVG     & 0.406 & 0.363 & 11.8\% \\
        \bottomrule
    \end{tabular}
    \caption{Performance metrics and calibration loss (calculated as $\frac{\text{Actual} - \text{Min}}{\text{Min}}$) for the three models across different effective priors.}
    \label{tab:comparison_results}
\end{table}

\section{Evaluation on Effective Priors}

\subsection*{Which models perform best?}

The intrinsic discriminative power of a model is measured by its Minimum DCF. The \textbf{Standard MVG} and the \textbf{Naive Bayes} classifiers are unambiguously the best-performing models. At a uniform prior ($\tilde{\pi} = 0.5$), the Standard MVG achieves the lowest overall minimum cost at $0.130$, almost identically matched by Naive Bayes at $0.131$. At the extreme prior of $\tilde{\pi} = 0.1$, Naive Bayes takes the lead with a Minimum DCF of $0.257$ compared to the Standard MVG's $0.263$. 

The success of the Naive Bayes classifier is theoretically supported by the correlation analysis conducted in previous evaluations, which revealed that the maximum off-diagonal correlation coefficient in the dataset is a negligible $0.0489$. Because the features are inherently uncorrelated, the conditional independence assumption of the Naive Bayes model holds perfectly, allowing it to perform optimally without overfitting cross-feature noise.

Conversely, the \textbf{Tied MVG} is the worst-performing model across all scenarios, peaking at a Minimum DCF of $0.442$ at $\tilde{\pi} = 0.9$. This failure is rooted in the empirical variances of the data: Genuine Feature 1 has a variance of $1.4302$ while Fake Feature 1 has a variance of just $0.5696$. The Tied MVG enforces a strictly shared within-class covariance matrix, destroying the ability to model these distinct class spreads and resulting in a sub-optimal linear boundary.

\subsection*{Are the relative results of the models consistent over the different applications?}

Yes, the relative performance rankings are highly consistent across the different target applications. The Standard MVG and Naive Bayes classifiers consistently form the top tier of performance, always exhibiting Minimum DCFs within $0.01$ of each other regardless of the effective prior. The Tied MVG classifier is equally consistent in its inferiority, remaining the worst-performing model across all three effective prior configurations by a significant margin. 

\subsection*{Now consider also actual DCFs. Are the models well calibrated (i.e., with a calibration loss in the range of few percents of the minimum DCF value) for the given applications?}

The models are well-calibrated \textit{only} at the uniform effective prior ($\tilde{\pi} = 0.5$). At this operational point, the Actual DCF closely tracks the Minimum DCF, resulting in minor calibration losses of 7.7\% for the Standard MVG, 9.9\% for Naive Bayes, and a mere 2.7\% for the Tied MVG. 

However, the models suffer a severe loss of calibration at the extreme priors ($\tilde{\pi} = 0.9$ and $\tilde{\pi} = 0.1$). For instance, the Standard MVG at $\tilde{\pi} = 0.9$ yields an Actual DCF of $0.400$ against a Minimum DCF of $0.342$, representing a massive 17.0\% calibration loss. This indicates that the theoretical decision threshold $t = -\log \frac{\tilde{\pi}}{1 - \tilde{\pi}}$ applied directly to the raw log-likelihood ratios becomes highly sub-optimal when the application is heavily skewed.

\subsection*{Are there models that are better calibrated than others for the different applications?}

Yes, there is a clear inverse relationship between a model's discriminative power and its calibration stability. The worst-performing model, the \textbf{Tied MVG}, is consistently the best-calibrated model across the applications. Its calibration loss remains remarkably low at $\tilde{\pi} = 0.5$ (2.7\%) and $\tilde{\pi} = 0.9$ (4.7\%), and it exhibits the lowest calibration loss at the $\tilde{\pi} = 0.1$ extreme (11.8\%). 

Because the Tied MVG uses a simplified, linear decision boundary derived from a pooled covariance matrix, its output scores are inherently more rigid and less susceptible to the extreme probabilistic shifts that degrade the calibration of the more complex, quadratic boundary of the Standard MVG.

\section{Extended Range Analysis (Bayes Error Plots)}

To fully contextualize the calibration gaps and performance consistency, we compute the Bayes Error plots over an extended range of prior log-odds from $-4$ to $+4$. These plots illustrate the continuous trajectory of the Actual DCF (red) and Minimum DCF (blue).

\begin{figure}
    \centering
    \includegraphics[width=0.75\textwidth]{Bayes_Error_Standard_MVG.png}
    \caption{Bayes Error plot for the Standard MVG classifier. Note the increasing deviation of the red Actual DCF curve from the optimal blue Minimum DCF curve at extreme log-odds.}
    \label{fig:bayes_mvg}
\end{figure}

\begin{figure}
    \centering
    \includegraphics[width=0.75\textwidth]{Bayes_Error_Naive_Bayes.png}
    \caption{Bayes Error plot for the Naive Bayes classifier. The performance profile is virtually identical to the Standard MVG, confirming that the independence assumption holds perfectly.}
    \label{fig:bayes_nb}
\end{figure}

\begin{figure}
    \centering
    \includegraphics[width=0.75\textwidth]{Bayes_Error_Tied_MVG.png}
    \caption{Bayes Error plot for the Tied MVG classifier. While the overall cost is higher, the red and blue curves remain closer together compared to the unconstrained models, indicating a higher robustness in calibration.}
    \label{fig:bayes_tied}
\end{figure}

\subsection*{What do you observe?}

Across all three classifiers, the minimum achievable cost sits at the center of the spectrum (prior log-odds $= 0$, corresponding to a uniform effective prior of $\tilde{\pi}=0.5$), yielding the characteristic "V" shape of the Minimum DCF curve. As the operational applications become more extreme (moving towards $-4$ or $+4$), both the theoretical Minimum DCF and the Actual DCF increase. The most striking observation is the diverging geometry of the curves: while the Actual and Minimum DCF overlap cleanly at the center, the red Actual DCF curve begins to pull away and lift higher than the blue minimum curve as we approach the extreme log-odds.

\subsection*{Are model rankings consistent across applications (minimum DCF)?}

Yes, visual inspection of the blue Minimum DCF curves across the entire log-odds range completely validates the consistency of the model rankings. The Standard MVG (Figure \ref{fig:bayes_mvg}) and Naive Bayes (Figure \ref{fig:bayes_nb}) blue curves trace nearly identical, low-profile paths across the entire spectrum from $-4$ to $+4$. When compared to the Tied MVG (Figure \ref{fig:bayes_tied}), the blue curve of the Tied model is visibly elevated at all points along the x-axis, confirming that the Tied MVG's inability to model class-specific variances universally harms its discriminative capacity regardless of the target application.

\subsection*{Are models well-calibrated over the considered range of applications?}

No, the extended Bayes error plots definitively prove that the models are not well-calibrated over the full range of applications. Calibration loss is visually represented by the vertical gap between the red (Actual) and blue (Minimum) curves. While this gap is zero near log-odds $0$, it grows exponentially towards the edges. 

This is most dramatically visible in the Tied MVG classifier (Figure \ref{fig:bayes_tied}): at a prior log-odds of $+4$, the Actual DCF reaches a maximum cost of $1.0$, while the theoretical Minimum DCF curve sits substantially lower at approximately $0.8$. The Standard MVG and Naive Bayes models also exhibit this widening trumpet-shape at the extremes. This geometric divergence visually proves that for highly asymmetric costs or severe class imbalances, the native log-likelihood ratios generated by these Gaussian models are no longer strictly aligned with the theoretical Bayes threshold, and a separate score re-calibration step would be required to deploy them effectively in a real-world edge case.